\documentclass{article}
\usepackage{amsmath}
\usepackage{amssymb}
\usepackage{graphicx}

\title{\textbf{Example Physics Research Project}}
\author{}
\date{}

\begin{document}

\title{\textbf{Example Physics Research Project}}

\section*{Overview}
This is an example project template. Replace this with your own research content. The system automatically extracts the Overview section as a summary on the project list.

Add images, equations, and formatted text to showcase your research. All LaTeX formatting will be automatically rendered on the website.

\section{Introduction}

Welcome to your research portfolio! This template shows you how to structure your research projects for the Event Horizon portfolio website.

The portfolio automatically parses LaTeX documents and renders them beautifully with:
\begin{itemize}
    \item Mathematical equation rendering
    \item Image inclusion and display
    \item Professional typography
    \item Responsive design for all devices
\end{itemize}

\section{Mathematics Example}

Here's an example of inline mathematics: $E = mc^2$, which can appear anywhere in your text.

Display mathematics is shown centered and prominent:

\begin{equation}
    F = ma
\end{equation}

You can also use alignment:

\begin{align}
    a &= \frac{v - v_0}{t} \\
    v &= v_0 + at \\
    s &= v_0 t + \frac{1}{2}at^2
\end{align}

\section{Equations with References}

\begin{equation}
    \label{eq:einstein}
    E^2 = (pc)^2 + (mc^2)^2
\end{equation}

You can reference equations like Equation \ref{eq:einstein} in your text.

\section{Project Structure}

Your project folder should contain:

\begin{enumerate}
    \item \textbf{Project_Name.tex} - Your main LaTeX file (required)
    \item \textbf{image1.png, image2.jpg} - Images referenced in the document
    \item \textbf{subdirectory/} - Optional folder for organizing images
\end{enumerate}

\section{Text Formatting}

You can use various text formatting options:
\begin{itemize}
    \item \textbf{Bold text} for emphasis
    \item \textit{Italic text} for different emphasis
    \item \texttt{Monospace} for code or technical terms
\end{itemize}

\section{Images}

Include images using the standard LaTeX command:

\begin{figure}
    \centering
    \includegraphics{placeholder.png}
    \caption{Example image caption - replace with your actual image}
\end{figure}

\section{Tips for Success}

\begin{enumerate}
    \item \textbf{Keep sections organized} - Use \texttt{\textbackslash section} and \texttt{\textbackslash subsection}
    \item \textbf{Add descriptive captions} - Figures should be self-explanatory
    \item \textbf{Use meaningful labels} - Reference equations and figures by label
    \item \textbf{Optimize images} - Compress before adding to project folder
    \item \textbf{Use standard LaTeX} - Avoid custom packages if possible
\end{enumerate}

\section{Supported Features}

\subsection{Equations}

All major equation environments are supported:
\begin{align*}
    \text{equation, align, gather, eqnarray, displaystyle}
\end{align*}

\subsection{Lists}

Both itemize and enumerate lists work perfectly with proper styling.

\subsection{Citations and References}

Use standard citation commands:
\begin{itemize}
    \item \verb|\cite{ref}| for general citations
    \item \verb|\citep{ref}| for parenthetical citations  
    \item \verb|\citet{ref}| for textual citations
\end{itemize}

\subsection{Figures}

You can include figures with:
\begin{itemize}
    \item PNG, JPG, JPEG, GIF, SVG, WebP formats
    \item Automatic scaling and optimization
    \item Captions and figure references
\end{itemize}

\section{Next Steps}

\begin{enumerate}
    \item \textbf{Customize} - Replace this template with your actual research
    \item \textbf{Add images} - Place your figures in the project folder
    \item \textbf{Build locally} - Run \verb|npm run dev| to preview
    \item \textbf{Deploy} - Push to GitHub or run \verb|npm run deploy|
\end{enumerate}

\section{Questions or Issues?}

\begin{itemize}
    \item Check SETUP\_GUIDE.md for detailed setup instructions
    \item Review README.md for feature documentation
    \item Ensure your .tex file syntax is correct
    \item Verify images exist in the project folder
\end{itemize}

\end{document}
